\section{Pendahuluan}

\subsection{Latar belakang dan tujuan}
Pengolahan citra digital memungkinkan perubahan intensitas, pemisahan objek,
dan pengukuran bentuk dilakukan secara komputasional. Laporan ini membahas
penyelesaian tiga bagian Ujian Tengah Semester mata kuliah Pengolahan Citra
Digital, yaitu selisih respons \textit{Laplacian of Gaussian} pada soal nomor 1,
segmentasi karakter pada soal nomor 2, serta penghitungan koin menurut ukuran pada soal nomor 3 \cite{soal-uts}.

Tujuan pekerjaan adalah mengimplementasikan algoritma dalam Python,
menjelaskan tahapan dan dasar matematisnya, serta memeriksa hasil dengan data
contoh dan pengujian otomatis. Pembahasan nomor 1 mempertahankan standar
deviasi yang tertulis pada soal, yaitu $0{,}1$ dan $0{,}2$. Pembahasan nomor 3
menggunakan fungsi \kode{regionprops} untuk mengukur setiap objek koin yang
telah dipisahkan.

\subsection{Ruang lingkup dan data}
Laporan ini mencakup nomor 1, 2, dan 3. Kode, citra masukan, serta hasil
perhitungan tersimpan dalam direktori masing-masing soal. Nomor 1 dan 3 dikerjakan oleh Bambang Herlambang (5803024019),
sedangkan nomor 2 dikerjakan oleh Dustin Pranata Wangsawidjaja (5803024021).

\begin{table}[H]
\centering
\caption{Bahan masukan dan implementasi yang digunakan.}
\label{tab:bahan}
\small
\begin{tabularx}{\textwidth}{@{}p{2.1cm}X@{}}
\toprule
\textbf{Bagian} & \textbf{Bahan dan program}\\
\midrule
Nomor 1 &
Citra sintetis grayscale $384 \times 256$ piksel pada
\path{soal1/data/citra_uji.png}; program \path{soal1/main.py}.\\
\addlinespace
Nomor 2 & Citra teks pada \path{soal2/teks.png}; program
\path{soal2/main.py}, dikerjakan oleh Dustin.\\
\addlinespace
Nomor 3 &
Citra sintetis $1100 \times 1100$ piksel berisi 100 koin terpisah pada
\path{soal3/data/koin_100.png}; program \path{soal3/main.py}.\\
\addlinespace
Validasi &
Metadata koin pada \path{soal3/data/koin_100_ground_truth.json} serta
pengujian pada \path{soal1/tests/} dan \path{soal3/tests/}.\\
\bottomrule
\end{tabularx}
\end{table}

Citra contoh nomor 1 dan 3 dibuat secara sintetis melalui \path{buat_data_contoh.py}
dengan \textit{seed} 9 agar dapat direproduksi. Citra koin bukan foto uang
logam nyata. Data acuan atau \textit{ground truth} hanya dipakai untuk
memeriksa hasil; program penghitung membaca citra tanpa mengambil jumlah
atau ukuran koin dari metadata acuan.

Kode lengkap dan berkas pendukung tersedia pada
\href{https://github.com/Bambang-code/uts-pcd-kelompok-9}{repositori GitHub Kelompok 9}.
Langkah pemasangan dan pelaksanaan disajikan pada
Bagian~\ref{sec:reproduksi}.
